% TOP_PROBLEM: {"id": 1465, "title": "O-Minimal Theory with Trans-Exponential Growth", "category_id": 4, "category": {"id": 4, "name": "algebra", "display_name": "Algebra", "description": "Group theory, ring theory, field theory, and algebraic structures.", "slug": "algebra", "order_index": 4, "created_at": "2026-07-31T15:26:25.671Z"}, "status": "open"}
% FIRST_POSTED: 2026-10-01
% ATTEMPT_STATUS: unresolved
% Imported from: unsolvedmath_additional_500.tex; UnsolvedMath ID 1465
% !TeX program = lualatex
% Compile twice: lualatex 1465.tex
% Self-contained: no external bibliography, images, or \input files.
% Numeric section labels and visible numbers are original UnsolvedMath IDs.
\documentclass[11pt,a4paper]{article}
\usepackage[margin=24mm,headheight=14pt]{geometry}
\usepackage{fontspec}
\setmainfont{Latin Modern Roman}
\setsansfont{Latin Modern Sans}
\setmonofont{Latin Modern Mono}
\usepackage{amsmath,amssymb,mathtools}
\usepackage{unicode-math}
\setmathfont{Latin Modern Math}
\usepackage{microtype}
\usepackage{enumitem,xurl,needspace}
\usepackage[dvipsnames]{xcolor}
\usepackage{hyperref}
\hypersetup{unicode=true,colorlinks=true,linkcolor=MidnightBlue,urlcolor=MidnightBlue,
 pdftitle={UnsolvedMath Problem 1465},pdfauthor={Source-annotated compilation},bookmarksnumbered=true}
\usepackage{bookmark,tocloft,titlesec,fancyhdr}
\setlength{\cftsecnumwidth}{4.2em}
\cftsetpnumwidth{2.5em}
\cftsetrmarg{4em}
\titleformat{\section}{\Large\bfseries\raggedright}{\thesection}{0.7em}{}
\pagestyle{fancy}\fancyhf{}
\fancyhead[L]{\small\sffamily Additional UnsolvedMath statements}
\fancyhead[R]{\small Original catalogue IDs}
\fancyfoot[C]{\thepage}
\setlength{\parindent}{0pt}
\setlength{\parskip}{5pt plus 1pt}
\setlength{\emergencystretch}{3em}
\setcounter{tocdepth}{1}\setcounter{secnumdepth}{1}
\setlist[itemize]{leftmargin=3em,itemsep=4pt,topsep=4pt,parsep=0pt}
\allowdisplaybreaks
\newcommand{\N}{\mathbb{N}}
\newcommand{\Z}{\mathbb{Z}}
\newcommand{\Q}{\mathbb{Q}}
\newcommand{\R}{\mathbb{R}}
\newcommand{\C}{\mathbb{C}}
\newcommand{\F}{\mathbb{F}}
\newcommand{\A}{\mathbb{A}}
\newcommand{\PP}{\mathbb{P}}
\newcommand{\E}{\mathbb{E}}
\newcommand{\eps}{\varepsilon}
\newcommand{\dd}{\,\mathrm d}
\newcommand{\sref}[2]{\hyperlink{source:#1:#2}{[S#2]}}
\newif\ifshowcatalogue
\showcataloguetrue
\newcommand{\cataloguescope}[1]{\ifshowcatalogue
 \paragraph{Statement recorded in the supplied source.}
 #1\fi}
\begin{document}
% UnsolvedMath ID 1465; source catalogue code MOD-008

\setcounter{section}{1464}

\section{An O-Minimal Trans-Exponential Expansion}\label{1465}

{\small\sffamily UnsolvedMath ID 1465 · Algebra}

\subsection{Definitions and mathematical statement}

\paragraph{Shared notation.}Unless a section says otherwise, $\N=\{1,2,\ldots\}$,
$\Z$ is the ring of integers, $\Q,\R,\C$ are the rational, real, and complex
fields, $[N]=\{1,\ldots,\lfloor N\rfloor\}$, and $\log$ is the natural
logarithm. For real functions of a parameter tending to infinity,
$f\ll g$ and $f=O(g)$ mean $|f|\leq Cg$ eventually for some fixed $C$;
$f\gg g$ reverses this comparison; $f=o(g)$ means $f/g\to0$;
$f\sim g$ means $f/g\to1$; and $f\asymp g$ means both order bounds.
A subscript on $O$ or $\ll$ specifies allowed constant dependence.
The expression $n^{o(1)}$ in an upper bound means growth smaller than
$n^\epsilon$ for each fixed $\epsilon>0$. Local definitions govern any
exception, finite-field convention, or use of the same symbol for another
object. An asymptotic claim is not automatically an effective algorithm.



An expansion of the ordered real field is o-minimal if each definable subset of the line, allowing parameters, is a finite union of points and intervals. A trans-exponential function $f$ eventually dominates every finite iterate of the usual exponential: for each $k\geq1$, $f(x)>\exp^{\circ k}(x)$ for all sufficiently large $x$. Seek an o-minimal expansion defining such a function. \sref{1465}{1} \sref{1465}{2}

\paragraph{Statement recorded in the supplied source.}
Does there exist an o-minimal first-order theory with a trans-exponential (rapid growth) function? \sref{1465}{1}

\subsection{Short English statement}

Can an o-minimal expansion of the real field define a function that eventually grows faster than every fixed-height exponential tower?

\subsection{Research attempt: a trans-exponential candidate whose derivative singularities define the integers}
\textbf{Status and source scope.} We construct a concrete rapidly growing function and prove its expansion cannot be o-minimal. This refutes a natural construction strategy, not the existence asked for. The inspected revised abstract of Fu [S2] reduces o-minimality of an analytic candidate to a regular-value condition; it does not announce unconditional o-minimality.

\subsubsection{1. A continuous increasing tower interpolation}
Fix $a\geq1$, set $c=e^a-a>0$, and let $h(u)=a+cu$ for $0\leq u\leq1$. For $x=n+u$, $n\in\mathbb N_0$ and $0\leq u<1$, put
\[f(x)=\exp^{\circ n}(h(u)).\tag{1}\]
At an integer endpoint the left limit is $\exp^{\circ(n-1)}(e^a)=\exp^{\circ n}(a)$, the value from the right. Thus $f$ is continuous and strictly increasing on $[0,\infty)$, and $f(x+1)=\exp(f(x))$.

For a fixed $k$ and $n\geq k$, applying $k$ logarithms to (1) gives at least $\exp^{\circ(n-k)}(a)$. For $t\geq1$, $e^t\geq2t$, so this quantity is at least $2^{n-k}a$. For sufficiently large $n$ it exceeds $n+1\geq x$. Applying $k$ exponentials back proves $f(x)>\exp^{\circ k}(x)$ eventually. Thus the growth requirement itself has been verified.

\subsubsection{2. Every positive integer leaves a definable defect}
Write $E_n=\exp^{\circ n}$ and $E_0$ for the identity. On each open unit interval $f$ is analytic. At the positive integer $n$, the one-sided derivatives are
\[f'_-(n)=cE'_{n-1}(e^a),\qquad
f'_+(n)=cE'_n(a)=ce^a E'_{n-1}(e^a).
\]
They are positive and unequal, with ratio $e^a>1$. Hence the nondifferentiability set of $f$ in $(0,\infty)$ is exactly $\{1,2,3,\ldots\}$.

Differentiability is first-order definable from the field operations, order and graph of $f$: there exists a real $L$ such that for every $\epsilon>0$ there is $\delta>0$ with
\[0<|t-x|<\delta\ \Longrightarrow
\left|\frac{f(t)-f(x)}{t-x}-L\right|<\epsilon.
\]
All quantifiers range over reals. Consequently the expansion by $f$ defines its nondifferentiability set, an infinite discrete subset of the line. Such a set is neither a finite union of points nor a union containing an interval. The expansion is therefore not o-minimal. Defining $f$ arbitrarily on the negative half-line cannot remove this already definable obstruction on $(0,\infty)$.

\subsubsection{3. What a smoother interpolation would have to change}
Replacing the linear $h$ by a differentiable connector with $h(0)=a$ and $h(1)=e^a$ removes the first-derivative jump only if
\[h'(1)=e^a h'(0).
\]
Higher smoothness imposes further endpoint identities obtained by differentiating $f(x+1)=e^{f(x)}$. Even satisfying all such identities would establish smoothness, not o-minimality: other definable sets in the expansion would still need to be controlled. Conversely, the failure of the piecewise-linear connector says nothing about an analytic candidate whose integer joins are invisible.

\subsubsection{4. Verification and remaining gap}
The diagnostic checks the endpoint and derivative-ratio formulas for initial iterates and the elementary inequality $e^t\geq2t$ on the domain needed by the symbolic growth induction. The proof does not rely on computing enormous tower values. It distinguishes the two goals explicitly: (1) has the required domination, but its singularity set violates o-minimality. An o-minimal construction must avoid this definable discrete set and every other such obstruction; proving that for any candidate remains absent. The result here is a rigorous failure analysis of one natural tower interpolation, not a nonexistence theorem for all trans-exponential expansions.

\subsection{Sources}

{\small

\begin{itemize}

\item[\hypertarget{source:1465:1}{S1}] ULAM.AI, \emph{UnsolvedMath}, supplied \texttt{problems.json}, record 1465 (MOD-008), statement and background. The embedded literature-review date is 2026-08-17; that is the archive’s review date, not a fresh certification. Dataset landing page: \url{https://huggingface.co/datasets/ulamai/UnsolvedMath}.

\item[\hypertarget{source:1465:2}{S2}] Y. Fu, Towards Trans-Exponential O-minimal Expansion of the Real Field, arXiv:2604.03477 (2026). \url{https://arxiv.org/abs/2604.03477}. Bibliographic information imported from the supplied record.

\item[\hypertarget{source:1465:3}{S3}] Trans-exponential o-minimal structure status page. \url{https://www.emergentmind.com/open-problems/existence-of-transexponential-o-minimal-structure}. Bibliographic information imported from the supplied record.

\end{itemize}

}

\label{end:1465}
\end{document}
